% answers_short.tex —— 2007 年·普通高等学校招生考试·重庆卷·理科 · 简短答案速查
% 由 tools/make_short_answers.py 自 answers.tex 的 \gksolhead/\gkanswer 提取，
% 并与 2007/2007_数学真题_重庆理_答案要点.md 逐题比对后生成，勿手改（改源 MD 后重跑）。

\begin{gknotice}
\textbf{使用说明}：本节只给最终答案便于快速核对，完整推导见解析版。答案为编写者独立推导，\textbf{非官方标准答案}。
\end{gknotice}

\gkpart{选择题}
\begin{center}\begin{tabular}{ccccccccccc}
\toprule
题号 & 1 & 2 & 3 & 4 & 5 & 6 & 7 & 8 & 9 & 10 \\
\midrule
答案 & A & D & C & B & A & C & B & B & D & C \\
\bottomrule
\end{tabular}\end{center}

\gkpart{填空题}
\begin{itemize}[leftmargin=2.6em]
  \item[\textbf{11.}] $\dfrac45$。
  \item[\textbf{12.}] $7$。
  \item[\textbf{13.}] $[-1,0]$。
  \item[\textbf{14.}] $\dfrac{3+3\sqrt3}2$。
  \item[\textbf{15.}] $25$。
  \item[\textbf{16.}] $\dfrac{8\sqrt3}3$。
\end{itemize}

\gkpart{解答题}
\begin{itemize}[leftmargin=2.6em]
  \item[\textbf{17.}] （1）最大值 $3+2\sqrt3$，最小正周期 $\pi$；（2）$\tan\dfrac45\alpha=\sqrt3$。
  \item[\textbf{18.}] （1）$\dfrac3{11}$；（2）分布列见解析，$E\xi=\dfrac{29900}{11}$。
  \item[\textbf{19.}] （1）$\dfrac{2\sqrt{29}}{29}$；（2）$\dfrac{3\sqrt3}2$。
  \item[\textbf{20.}] （1）$a=12$，$b=-3$；（2）减区间 $(0,1)$，增区间 $(1,+\infty)$；（3）$c\leqslant-1$ 或 $c\geqslant\dfrac32$。
  \item[\textbf{21.}] （1）$a_n=3n-1$；（2）证明见解析。
  \item[\textbf{22.}] （1）$\dfrac{x^2}{36}+\dfrac{y^2}{27}=1$；（2）定值 $\dfrac23$。
\end{itemize}
